\section{Discussion}\label{sec:discussion}

\paragraph{Reconciling with Wu et al.~\cite{wu2025single}.}
Wu et al.~\cite{wu2025single} report that \GST{} outperforms matched sampling, with the same temperature, top-$p$ and top-$k$,
on three 32B reasoning models. % CITE(wu2025single)
The gain, averaged over their eight benchmarks, is $+0.42$ to $+1.05$ points, with no variance or significance % CITE(wu2025single)
reported; GPQA-Diamond shows the largest gains, up to $+5.05$ points (their Table~2). % CITE(wu2025single)
At 1.5B on GSM8K we find no detectable difference % SETUP
(accuracy comparison: 52 of 60 for both; coupled ablation: 88 vs 86 of 100, McNemar $p = 0.754$). % C16.1 C16.3 C23.1 C23.2
The two findings are compatible.
Wu et al.\ motivate Gumbel-Softmax through the Gumbel-max trick; its consequence is that the top token of \GST{} is
distributed exactly as a sampled token, which we confirm (flip rate 0.45 for both). % C22.4
Dirichlet noise has no such guarantee, and in our runs its top token flips 0.03 less often than a sampled token % C22.8
would.
Wu et al.\ also find Dirichlet no better than sampling, although our Dirichlet uses $\gamma = 1.0$ % SETUP
and theirs 4.0 by default. % CITE(wu2025single)
Any accuracy difference between \GST{} and sampling must therefore come from the residual.
Our coupled ablation holds the noise fixed and removes only the residual.
The model reads it, yet accuracy does not change detectably.
Our paired interval runs from 5 points below sampling to 9 points above it, a range that includes the gains % C23.17
Wu et al.\ report.
For the Gumbel variant this refines their single-threaded account: the top token is a sample, and the residual is
read but leaves accuracy unchanged within our power.
Three differences could separate the settings.
A 32B model may use a residual that % CITE(wu2025single)
a 1.5B model only stores. % SETUP
GSM8K is near ceiling for our model (86 to 88\% in the coupled ablation), while the largest reported gains are on harder % C23.1
benchmarks.
Our study lacks the power to resolve effects of a few points, and Wu et al.\ report no variance.
Running the coupled ablation on a 32B model and a harder benchmark would show which explanation applies. % CITE(wu2025single)

\paragraph{The tension with Rizvi-Martel et al.~\cite{rizvimartel2026illusion}.}
They conclude that off-the-shelf models collapse \ST{} inputs within the first few layers, while our
\cref{fig:snap} shows a gradual collapse that is complete only at the output.
The two results measure different things.
Their evidence for early collapse comes from vanilla \ST{}, every 50 decoding steps: the logit-lens entropy of % CITE(rizvimartel2026illusion)
separate soft and discrete runs, and the KL divergence between forward passes with the soft input and with its top
token.
Those inputs are close to one token on average (mean mixing-weight entropy 0.10 to 0.18 nats), so little blend is % CITE(rizvimartel2026illusion)
left to lose.
We condition on forks, feed exact 50/50 blends of the two candidates, and track the position between them at every
layer.
On these inputs the blend is still far from a clean switch at layer 7 (snap 0.37 to 0.70). % C21.4
This agrees with their own test on uniform mixtures, where entropy drops only at the last layers, and with
Backour~\cite{backour2026dynamics}, who finds that Qwen models accumulate collapse toward the output.
Collapse in the next-token readout also does not mean the runner-up is lost: a trace of it persists, mostly in the
keys and values of layers 0 to 7, and shifts later predictions toward it (\cref{sec:q2}). % C21.10
Our blends are 50/50, likely an upper bound on real forks, where the median residual is 0.12. % C22.12

\paragraph{What our evidence says about superposition.}
Our results separate three levels of token-level superposition: a second candidate present in the input, read by
the model, and used.
The forced 50/50 blends are present in the input, and a trace of the runner-up persists in the cache.
The residual of \GST{} is read: it shifts next-token predictions far more than a random residual.
We find no evidence that it is used.
For that, the trace would have to shift predictions on the model's own path, where we detect no shift, and the
residual would have to change answers, where we detect no accuracy change and the real residual changes answers
no more often than a random one.
Two limits keep this from being a general verdict on superposition.
Our blends are forced, not learned; models trained with latent thoughts from scratch show signs of using
superposition \cite{rizvimartel2026illusion}, and reinforcement learning with soft tokens is now practical
\cite{butt2025soft,zheng2025softgrpo}.
And we test single positions, not whole reasoning paths: a model could in principle keep two lines of reasoning
alive across many steps without either one leaving a detectable trace at a single fork.
The coupled ablation applies to trained models unchanged: shared noise gives the arms the same top token until their
contexts diverge, so the residual is the only source of divergence.
